\section{Domain Task 9: Phân tích danh mục chủ nhà (Host Portfolio Analysis)}

Câu hỏi nghiệp vụ: Nhóm chủ nhà cá nhân và chủ nhà chuyên nghiệp có cơ
cấu phân bố như thế nào, và có sự chênh lệch ra sao về hiệu suất kinh
doanh (điểm đánh giá, tỷ lệ lấp đầy)?

Mục đích: Đánh giá mức độ thương mại hóa của thị trường Airbnb NYC và so
sánh lợi thế cạnh tranh giữa mô hình kinh doanh cá thể so với tổ
chức/chuỗi phòng.

\subsection{Biểu đồ 2: So sánh hiệu suất kinh doanh theo loại chủ nhà (Bar Chart)}

\subsubsection*{A. Idiom}

\begin{longtable}{@{} p{2.8cm} p{12cm} @{}}
\toprule
\textbf{Đặc điểm} & \textbf{Chi tiết} \\
\midrule
\endhead

\bottomrule
\endfoot

\textbf{Idiom} & Bar Chart \\
\addlinespace[0.2cm]

\textbf{What} & Loại chủ nhà (Host Type): C, Điểm đánh giá (Avg Rating): Q, Tỷ lệ lấp đầy (Avg Is Booked): Q \\
\addlinespace[0.2cm]

\textbf{Why} & produce (derive) $\rightarrow$ lookup $\rightarrow$ compare
\begin{itemize}[nosep, leftmargin=*, topsep=0.1cm]
    \item derive: tính trung bình cho Rating và tỉ lệ lấp đầy
    \item Compare: trực diện chất lượng dịch vụ và khả năng thu hút khách hàng giữa 2 nhóm chủ nhà
\end{itemize} \\
\addlinespace[0.2cm]

\textbf{How} & \textbf{Encode:}
\begin{itemize}[nosep, leftmargin=*, topsep=0.1cm]
    \item Mark: Line (Bar mark)
    \item Channel:
    \begin{itemize}[nosep, leftmargin=0.5cm]
        \item PosX: Phân biệt 2 loại chủ nhà
        \item PosY: thể hiện độ lớn của Avg Rating / Avg Is Booked
        \item Hue Color: Phân biệt loại chủ nhà
    \end{itemize}
\end{itemize}

\vspace{0.2cm}
\textbf{Manipulate:}
\begin{itemize}[nosep, leftmargin=*, topsep=0.1cm]
    \item Hover + Tooltip: Xem chi tiết các con số
    \item Facet (Juxtapose): tách thành 2 biểu đồ xếp dọc nhau (1 cho rating, 1 cho IsBooked) dùng chung trục X để dễ đối chiếu
\end{itemize} \\
\addlinespace[0.2cm]

\textbf{Scale} &
\begin{itemize}[nosep, leftmargin=*, topsep=0pt]
    \item Main key: 2 (Loại chủ nhà)
\end{itemize} \\
\end{longtable}

\begin{figure}[H]
    \centering
    \safeincludegraphics[width=0.9\linewidth]{charts/domain_task9_chart2.png}
    \caption{Hình 9.2. Biểu đồ so sánh hiệu suất kinh doanh theo loại chủ nhà}
    \label{fig:domain-task9-chart2}
\end{figure}

\subsubsection*{B. Đánh giá biểu đồ}

\textbf{1. Nguyên lý biểu đạt:}

\begin{itemize}
    \item Thuộc tính: Loại chủ nhà (C), các chỉ số hiệu suất (Q).
    \item Channel:
    \begin{itemize}
        \item PosX: Định vị nhóm $\rightarrow$ dùng cho thuộc tính (C).
        \item PosY: Chiều dài cột $\rightarrow$ dùng cho thuộc tính định lượng (Q).
        \item Color (Hue): Nhấn mạnh sự phân loại $\rightarrow$ dùng cho thuộc tính (C).
    \end{itemize}
    \item Đánh giá mức độ quan trọng:
    \begin{itemize}
        \item Ưu tiên 1 -- Chiều dài (PosY / Length): Hiệu suất kinh doanh (Q) là thông tin cốt lõi nhất cần so sánh. Theo Mackinlay, Position/Length trên một trục chung là kênh mạnh nhất, chính xác tuyệt đối nhất cho dữ liệu (Q).
        \item Ưu tiên 2 -- Vị trí không gian (PosX): Dùng để chia tách 2 loại chủ nhà (C). Kênh Position cũng là kênh mạnh nhất cho dữ liệu định danh (Categorical).
        \item Ưu tiên 3 -- Màu sắc (Hue Color): Việc dùng Hue để tô màu cho các cột là một dạng mã hóa lặp (Redundant Encoding). Tuy không bắt buộc do đã có PosX phân tách, nhưng nó đóng vai trò liên kết thị giác (nhấn mạnh Xanh = Cá nhân, Cam = Chuyên nghiệp) xuyên suốt từ biểu đồ 9.1 sang.
    \end{itemize}
\end{itemize}

\textbf{2. Nguyên lý hiệu quả:}

\begin{itemize}
    \item Độ chính xác (Accuracy): Trục PosY biểu diễn thông qua chiều dài (Length) bám sát gốc 0. Theo định luật Stevens, Length có $n = 1.0$, do đó mắt người cảm nhận cực kỳ chính xác. Độ lỗi Log\_error rất thấp.
    \item Khả năng phân biệt (Discriminability): 2 cột đứng cạnh nhau với 2 màu tương phản mạnh (Xanh/Cam) giúp mắt người đọc ngay lập tức nhận ra bên nào cao hơn mà không cần suy nghĩ.
    \item Khả năng tách biệt (Separability): Sử dụng Facet (Juxtapose) giải quyết việc tách biệt 2 hệ đo lường khác nhau (Rating thang 5, IsBooked là tỷ lệ \%) $\Rightarrow$ không xảy ra hiện tượng nhiễu trục.
\end{itemize}

\subsubsection*{C. Phân tích biểu đồ (Insight)}

\begin{itemize}
    \item Dù yếu thế về số lượng nguồn cung, chủ nhà ``Cá nhân'' (màu xanh) lại mang đến trải nghiệm tốt hơn so với nhóm chuyên nghiệp.
    \item Hệ quả là tỷ lệ lấp đầy của chủ nhà ``Cá nhân'' cũng cao hơn nhóm chuyên nghiệp. Điều này chứng minh khách hàng đề cao tính cá nhân hóa, sự thân thiện của các cá thể hơn là dịch vụ rập khuôn của các tổ chức.
\end{itemize}
